\section{全球具身智能产业链 Mapping}

本章整理谢晨给出的产业链图谱。具身智能之所以会拆成多类公司，是因为它同时要发明“下一代车的平台”和“这个平台上的 L4/L5 算法”。自动驾驶至少有车这个成熟平台；具身智能连平台、本体、传感器、手、模型、数据和场景都在早期。难度过高，就会推动产业链分工。

节目里大致分出四类公司：硬件公司、Foundation Model 公司、垂直场景软硬结合公司、以仿真为中心的基础设施公司。这个分类不是静态标签，而是说明不同公司把风险放在哪里：硬件公司解决本体和供应链，模型公司解决大脑和数据 recipe，落地公司解决场景 ROI，仿真公司解决数据和训练世界。

\begin{figure}[H]
\centering
\includegraphics[width=0.96\textwidth]{embodied-supply-chain-map.png}
\caption{全球具身智能产业链：硬件、基座模型、垂直落地和仿真公司分工。自制概念图，依据 00:55:25--01:09:22 对谈内容整理。}
\end{figure}

\begin{knowledgebox}{读图：四类公司对应四种稀缺资源}
硬件层稀缺的是可靠本体、成本和供应链；模型层稀缺的是算力、人才和通用数据 recipe；垂直落地层稀缺的是客户、场景和运营；仿真层稀缺的是高质量物理世界、工具链和评价闭环。具身智能太重，因此短期更像产业生态，而不是单家公司全包。
\end{knowledgebox}

\subsection{硬件、模型、落地与仿真}

本节把四类公司逐个展开。硬件公司代表是宇树和灵巧手公司，它们把机器人本体带给科研和产业界，某种程度上建立了研究标准。Foundation Model 公司包括 Physical Intelligence、Skild、NVIDIA/GEAR、Google/DeepMind 等，它们追求机器人“大脑”和 scaling recipe。软硬结合落地公司包括 Figure、Tesla Optimus、The Bot Company、DynaRobotics 等，重点是具体垂直场景。仿真公司如光轮和 Genesis，则围绕 Real2Sim、Simulation、Sim2Real 做基础设施。

讲者对宇树的评价很有产业启发：它面向国际学术界提供本体，让论文和学生都围绕宇树机器人形成习惯，再把这种习惯带入产业。这说明硬件层的“标准”不只来自参数，也来自开发者生态、论文生态和人才迁移。对应到仿真层，光轮想成为的也是这种标准：不是一个供应商，而是具身仿真的共同底座。

\begin{knowledgebox}{产业链角色表}
\begin{tabularx}{\textwidth}{p{0.18\textwidth}p{0.32\textwidth}X}
\toprule
角色 & 核心任务 & 主要竞争维度 \\
\midrule
硬件公司 & 机器人本体、手、传感器和采集设备 & 成本、可靠性、开发者生态、供应链。 \\
模型公司 & 机器人基座模型、大脑、VLA 和 scaling recipe & 算力、人才、数据金字塔、RL 能力。 \\
落地公司 & 面向主机厂、仓储、家庭、餐饮等场景交付 & 客户 ROI、软硬结合、运营和安全。 \\
仿真公司 & 生成可训练世界，支撑数据与算法闭环 & 物理资产、solver、API、Sim2Real 成功率。 \\
\bottomrule
\end{tabularx}
\end{knowledgebox}

\subsection{中美机会差异}

上一节是全球产业链，本节看国别结构。谢晨认为美国存在模型层创业机会，因为美国生态更强调产业分工，终端用户对软件、推理和自动化服务的付费能力更强。OpenAI、NVIDIA 这类公司能通过模型、算力和生态拿到很高收入，因此模型层可以成为独立公司。

中国的结构则不同。国内软件和 token 付费能力相对弱，硬件、供应链和整机能力更强，所以许多公司会走软硬结合、端到端交付和硬件销售路线。节目中提到字节、小米、理想等公司更可能适合做“大脑”，因为它们有算力、AI 人才、产品入口、供应链或车辆/机器人场景。这个判断不是说中国没有模型能力，而是说商业模式和产业组织会把模型能力放进更垂直的一体化体系。

\begin{figure}[H]
\centering
\includegraphics[width=0.94\textwidth]{china-us-embodied-opportunity.png}
\caption{中美具身机会差异：美国更有模型层机会，中国更适合大脑+本体+场景整合。自制概念图，依据 01:09:22--01:15:33 对谈内容整理。}
\end{figure}

\begin{knowledgebox}{读图：差异来自付费结构和产业分工}
左侧美国更容易出现独立模型层，因为软件付费、推理付费和生态分工支撑模型公司获得收入。右侧中国更容易把大脑、本体和场景整合，因为硬件供应链强、软件单独收费弱、产品公司更需要端到端价值。读图时不要把它理解成技术高低，而要理解成商业结构差异。
\end{knowledgebox}

\subsection{数据金字塔与殊途同归}

本节说明为什么不同技术路线可能最终收敛。表面上，有团队说自己走真实数据路线，有团队说自己走仿真路线，有团队强调 imitation learning，有团队强调 RL。但谢晨观察到，顶尖客户越来越在内部采用 data pyramid：底层是互联网或通用数据，中间是合成数据和仿真，顶部是少量高质量真实数据。

这个金字塔结构很符合前文逻辑。互联网数据提供世界知识和视觉/语言预训练，合成数据提供可控、多样、可并行的任务环境，真实数据提供 grounding、校准和最终验证。不同团队对 RL、VLA、合成比例和算法架构仍有差异，但“真实 + 合成 co-training”已经成为核心机构的共同认知。

\begin{importantbox}{数据金字塔}
具身智能可能不会纯靠真实数据，也不会纯靠仿真数据。更稳的路线是：底层通用数据建立世界知识，中层合成/仿真数据扩大任务分布，顶层真实数据做校准、示范和部署验证。
\end{importantbox}

\subsection{本章小结}

具身智能产业链会长期分工。硬件、模型、落地和仿真分别掌握不同稀缺资源；中美差异更多来自商业模式和产业组织；不同路线最终可能都收敛到数据金字塔和真实/合成混合训练。

